\section{Integrazioni e Modifiche all'Assignment 1}

A seguito dell'analisi delle risposte fornite dagli utilizzatori tramite i questionari di valutazione, sono emerse nuove esigenze che ci hanno spinto ad aggiornare i documenti redatti nel primo assignment. I feedback raccolti hanno evidenziato la forte necessità di espandere le funzionalità rivolte non solo al turista, ma anche a chi lavora quotidianamente all'interno del polo culturale ed è addetto alla sua gestione.

\subsection{Aggiunta del Task 4 per il gestore e relativi Sotto-Task}
Analizzando le necessità organizzative degli utenti, abbiamo deciso di ampliare le potenzialità del sistema introducendo un nuovo macro-task destinato specificamente al personale o al gestore della struttura (rispecchiando le esigenze della Persona \textbf{Giovanni}):

\begin{itemize}[nosep, leftmargin=1.4em]
    \item \textbf{T4: Gestione e authoring dei contenuti culturali.}
\end{itemize}

Per tradurre questo macro-task in azioni pratiche e valutabili all'interno dell'applicazione, lo abbiamo suddiviso nei seguenti sotto-task operativi:
\begin{itemize}[nosep, leftmargin=1.4em]
    \item \textbf{4.1 Creazione/modifica scheda opera}: la capacità di inserire e aggiornare agevolmente testi, immagini e coordinate per la Realtà Aumentata (AR) in totale autonomia.
    \item \textbf{4.2 Configurazione percorso AR guidato}: la possibilità di definire e riorganizzare visivamente l'ordine delle tappe e i punti interattivi, ad esempio per deviare il flusso dei visitatori lontano da eventuali aree in stato di manutenzione.
    \item \textbf{4.3 Monitoraggio statistiche di visita}: la consultazione dei dati sui flussi turistici per capire quali opere attraggono maggiormente l'attenzione, potendo così ottimizzare di conseguenza l'offerta culturale.
\end{itemize}

\subsection{Aggiornamento del form di Empowerment}
Avendo introdotto le nuove funzionalità di gestione previste dal Task 4, abbiamo riformulato di conseguenza anche il form dell'empowerment. 

In questo modo, intendiamo verificare se la nostra applicazione è in grado di conferire, anche a un impiegato privo di specifiche competenze informatiche, i "super-poteri" necessari per modellare percorsi, gestire informazioni digitali tridimensionali e intervenire per correggere tempestivamente inesattezze sui contenuti in esposizione.

\subsection{Aggiornamento del ruolo della Persona Giovanni}
A seguito di un'analisi più approfondita dei dati emersi dai questionari esplorativi somministrati alle Pubbliche Amministrazioni e agli addetti ai lavori, è stato necessario raffinare il profilo della Persona \textbf{Giovanni}. 

Originariamente delineato come un operatore di custodia e assistenza fisica nel museo, il ruolo di Giovanni è stato evoluto in quello di un \textbf{funzionario comunale} responsabile della digitalizzazione e della valorizzazione del territorio. Questa modifica, guidata dai feedback reali degli stakeholder, permette di inquadrare meglio le sfide gestionali e strutturali (come la scarsità di finanziamenti e la complessità dei processi burocratici) che Culturia mira a risolvere, posizionando l'applicazione come uno strumento di autonomia decisionale ed operativa per l'Ente Gestore.
